% sec5_unrolling.tex
% Section 5: Deep Unrolling
% Robust Sparse Convolution Coding for Radar Applications
% DSIPS2026 -- Tokyo, March 2026
%
% 8 slides:
%   1. What is Deep Unrolling?
%   2. From ISTA to LISTA: A Motivating Example
%   3. Unrolling for RPCA + Sparse Coding
%   4. LCRPCA: Unrolled Architecture
%   5. LCRPCA: Architecture Details  [skip?]
%   6. LCRPCA: Results with Scarce Data
%   7. Why Fewer Parameters Help
%   8. Comparison: Classical vs.\ Learned

\section{Deep Unrolling}

% ------------------------------------------------------------------
% SLIDE 1 -- What is Deep Unrolling?
% ------------------------------------------------------------------
\begin{frame}{What is Deep Unrolling?}
  \begin{columns}[T]
    % ------- LEFT: concept -------
    \column{0.52\textwidth}
      \vspace{0.3em}
      \textbf{Iterative optimization algorithm:}

      \vspace{0.4em}
      \begin{shibuiframe}{ }
        \[
          \mathbf{x}^{l+1}
            = \mathcal{A}\!\left(\mathbf{x}^{l};\,\boldsymbol{\theta}\right),
          \quad l = 1,\ldots
        \]
        \begin{itemize}
          \item $\boldsymbol{\theta}$: \emphtext{fixed} parameters
                (e.g.\ step size, threshold)
          \item Iteration count: \emphtext{variable} (run until convergence)
        \end{itemize}
      \end{shibuiframe}

      \vspace{0.6em}
      \textbf{Deep unrolling:}
      \begin{enumerate}
        \item \emphtext{Truncate} to $\nlayer$ iterations
        \pause
        \item Make parameters \emphtext{per-layer} and \emphtext{learnable}
      \end{enumerate}

      \vspace{0.4em}
      \begin{shibiunote}[Result]
        A neural network with $\nlayer$ layers,
        each implementing \emphtext{one algorithm step}.
      \end{shibiunote}

    % ------- RIGHT: properties -------
    \column{0.48\textwidth}
      \vspace{0.3em}
      \textbf{Advantages:}
      \begin{itemize}
        \item \emphtext{Interpretable}: physical meaning per layer
        \pause
        \item \emphtext{Data-efficient}: few parameters
        \pause
        \item \emphtext{Flexible}: learns from data
        \pause
        \item \emphtext{Fast}: fixed layers, no loops
      \end{itemize}

      \vspace{0.3em}
      \textbf{Limitations:}
      \begin{itemize}
        \item Fixed depth (no guarantee)
        \item Non-convex training
      \end{itemize}

      \vspace{0.4em}
      \citebadge{Gregor \& LeCun (2010) LISTA;
                 Monga et al.\ (2021) survey}
  \end{columns}
\end{frame}

% ------------------------------------------------------------------
% SLIDE 2 -- From ISTA to LISTA: A Motivating Example
% ------------------------------------------------------------------
\begin{frame}{From ISTA to LISTA: A Motivating Example}
  \begin{columns}[T]
    % ------- LEFT: ISTA -------
    \column{0.44\textwidth}
      \vspace{0.2em}
      \textbf{ISTA} solves:
      \[
        \min_{\mathbf{x}}
        \tfrac{1}{2}\ltwonorm{\mathbf{y}-\mathbf{W}\mathbf{x}}^{2}
        + \sparsityhp\lonenorm{\mathbf{x}}
      \]

      \vspace{0.1em}
      \begin{shibuiframe}{ISTA update}
        \[
          \mathbf{x}^{l+1}
            = \pxSoftT{\sparsityhp}\!\left\{
                \left(\mathbf{I}
                  - \tfrac{1}{\mu}\mathbf{W}^{\top}\mathbf{W}\right)
                \mathbf{x}^{l}
                + \tfrac{1}{\mu}\mathbf{W}^{\top}\mathbf{y}
              \right\}
        \]
      \end{shibuiframe}

      \vspace{0.3em}
      Fixed $\mathbf{W}$, $\mu$, $\sparsityhp$ -- iterate
      until convergence.

    % ------- ARROW -------
    \column{0.08\textwidth}
      \vspace{4em}
      \centering
      \large$\Rightarrow$\\[0.4em]
      \tiny Unroll\\+\\Learn

    % ------- RIGHT: LISTA -------
    \column{0.48\textwidth}
      \vspace{0.2em}
      \textbf{LISTA} (learnable ISTA):

      \vspace{0.1em}
      \begin{shibuiframe}{LISTA layer $\ilayer$}
        \[
          \mathbf{x}^{(\ilayer+1)}
            = \pxSoftT{\lparamlam{\ilayer}}\!\left(
                \mathbf{W}_{t}^{(\ilayer)}\mathbf{x}^{(\ilayer)}
                + \mathbf{W}_{e}^{(\ilayer)}\mathbf{y}
              \right)
        \]
        Learnable per layer:
        $\mathbf{W}_{t}^{(\ilayer)},\;
         \mathbf{W}_{e}^{(\ilayer)},\;
         \lparamlam{\ilayer}$
      \end{shibuiframe}

      \vspace{0.3em}
      \textbf{Training loss:}
      \[
        \lossfn
          = \tfrac{1}{N}\sum_{n}
            \ltwonorm{\hat{\mathbf{x}}^{n}
              - \mathbf{x}^{*n}}^{2}
      \]
  \end{columns}

  \vspace{0.1em}
  \begin{center}
  \small
  % unrolling_layers.tex
% Concept diagram: ISTA algorithm unrolled into T fixed layers.
% Top row: classical ISTA (loop with arbitrary iterations).
% Bottom row: LISTA / unrolled (fixed T layers, learnable params).
%
% Colors: theme names
% Bounding box: approx 11cm x 5.5cm

\begin{tikzpicture}[
    % Shared layer box style
    layerbox/.style = {
      draw=shibuiaccent,
      fill=shibuicodebg,
      rounded corners=3pt,
      thick,
      text=shibuitext,
      font=\scriptsize,
      align=center,
      minimum width=1.8cm,
      minimum height=1.1cm,
      inner sep=4pt
    },
    % Fixed-param box (lighter border)
    fixedbox/.style = {
      draw=shibuitext,
      fill=shibuicodebg,
      dashed,
      rounded corners=3pt,
      text=shibuitext,
      font=\scriptsize,
      align=center,
      minimum width=1.8cm,
      minimum height=1.1cm,
      inner sep=4pt
    },
    arr/.style = {->, draw=shibuitext, thick,
                  shorten >=2pt, shorten <=2pt},
    lbl/.style = {font=\small, text=shibuitext}
  ]

  % ========== TOP HALF: classical iterative algorithm ==========
  \node[lbl, font=\small\bfseries, text=shibuitext]
    at (0, 4.8) {Iterative algorithm (ISTA)};

  % Input
  \node[lbl, font=\scriptsize] (yin) at (-1.2, 3.7) {$\mathbf{y}$};

  % Three "iteration" boxes (same W each time)
  \foreach \i/\xi in {1/1.0, 2/3.2, 3/5.4} {
    \node[fixedbox] (ista\i) at (\xi, 3.7)
      {$\mathcal{S}_\lambda$\\[2pt]
       $\mathbf{W}^\top, \mathbf{W}$\\[1pt]
       fixed};
  }

  % Dots to indicate more iterations
  \node[lbl, font=\large] at (7.2, 3.7) {$\cdots$};

  % Output
  \node[lbl, font=\scriptsize] (xout_ista) at (8.6, 3.7) {$\hat{\mathbf{x}}$};

  % Arrows
  \draw[arr] (yin.east) -- (ista1.west);
  \draw[arr] (ista1.east) -- (ista2.west);
  \draw[arr] (ista2.east) -- (ista3.west);
  \draw[arr] (ista3.east) -- (7.0, 3.7);
  \draw[arr] (7.5, 3.7) -- (xout_ista.west);

  % Loop arrow (feedback)
  \draw[->, draw=shibuitext, thin, dashed, rounded corners=4pt]
    (5.4 + 0.9, 3.7 + 0.55) -- (5.4 + 0.9, 3.7 + 1.0)
    -- (1.0 - 0.9, 3.7 + 1.0)
    -- (1.0 - 0.9, 3.7 + 0.55);
  \node[lbl, font=\tiny, text=shibuitext]
    at (3.2, 3.7 + 1.15) {repeat $i = 1, \ldots$};

  % ========== UNROLLING ARROW ==========
  \draw[->, draw=shibuiaccent, very thick]
    (4.5, 2.95) -- (4.5, 2.3)
    node[lbl, right=4pt, midway, font=\small\bfseries,
         text=shibuiaccent] {unroll + learn};

  % ========== BOTTOM HALF: unrolled (LISTA / LCRPCA) ==========
  \node[lbl, font=\small\bfseries, text=shibuitext]
    at (0, 2.0) {Unrolled network (\nlayer{} layers)};

  % Input
  \node[lbl, font=\scriptsize] (yin2) at (-1.2, 1.1) {$\mathbf{y}$};

  % Three learnable layer boxes
  \foreach \i/\xi in {1/1.0, 2/3.2, 3/5.4} {
    \node[layerbox] (lista\i) at (\xi, 1.1)
      {$\mathcal{S}_{\lparamlam{\i}}$\\[2pt]
       $\Wone{\cdot}^{(\i)}$, $\Wtwo{\cdot}^{(\i)}$\\[1pt]
       \textbf{learned}};
  }

  % Dots
  \node[lbl, font=\large] at (7.2, 1.1) {$\cdots$};

  % Output
  \node[lbl, font=\scriptsize] (xout_lista) at (8.6, 1.1) {$\hat{\mathbf{x}}$};

  % Arrows
  \draw[arr] (yin2.east) -- (lista1.west);
  \draw[arr] (lista1.east) -- (lista2.west)
    node[lbl, font=\tiny, midway, below=2pt] {$\mathbf{x}^{(1)}$};
  \draw[arr] (lista2.east) -- (lista3.west)
    node[lbl, font=\tiny, midway, below=2pt] {$\mathbf{x}^{(2)}$};
  \draw[arr] (lista3.east) -- (7.0, 1.1);
  \draw[arr] (7.5, 1.1) -- (xout_lista.west);

  % Learnable parameters bracket under all layers
  \draw[decorate, decoration={brace, amplitude=5pt, mirror},
        shibuitext]
    (0.1, 0.45) -- (6.3, 0.45)
    node[lbl, font=\scriptsize, midway, below=6pt]
    {$\nparams$ learnable parameters
     $\{\lparamlam{t}, \lparammu{t}, \lparamw{t},
        \Wone{k}^{(t)}, \Wtwo{k}^{(t)}\}$};

\end{tikzpicture}

  \end{center}
\end{frame}

% ------------------------------------------------------------------
% SLIDE 3 -- Unrolling for RPCA + Sparse Coding
% ------------------------------------------------------------------
\begin{frame}{Unrolling for RPCA + Sparse Coding}
  \textbf{Context:} apply unrolling to the TWRI problem (KRPCA).

  \vspace{0.4em}
  \begin{shibuiframe}{PGD for KRPCA -- algorithm to unroll}
    \begin{align*}
      \Mwall^{(n+1)}
        &= \pxSvT{\lowrankhp}\!\left(
             \Mwall^{(n)}
             - \stepsize\bigl(
                 \Mtwri
                 - \Mwall^{(n)}
                 - \textstyle\sum_{\iatom}
                   \Mdict_{\iatom}\Fconv\Mtarget_{\iatom}^{(n)}
               \bigr)
           \right)
      \\[6pt]
      \Mtarget_{\iatom}^{(n+1)}
        &= \pxSoftT{\sparsityhp}\!\left(
             \Mtarget_{\iatom}^{(n)}
             - \stepsize
               \{\Mdict_{\iatom}\}^{\star}\!
               \bigl(
                 \Mtwri
                 - \Mwall^{(n)}
                 - \textstyle\sum_{\iatom}
                   \Mdict_{\iatom}\Fconv\Mtarget_{\iatom}^{(n)}
               \bigr)
           \right)
    \end{align*}
  \end{shibuiframe}

  \vspace{0.3em}
  where $\Mdict_{\iatom}$ is the physics-based propagation filter.

  \vspace{0.5em}
  \begin{columns}[T]
    \column{0.55\textwidth}
      \textbf{Key substitution -- replace fixed filters by
      learnable ones:}
      \[
        \stepsize\cdot\Mdict_{\iatom}
          \;\longrightarrow\;
          \Wone{\iatom}^{(\ilayer)},
        \qquad
        \stepsize\cdot\Mdict_{\iatom}^{\star}
          \;\longrightarrow\;
          \Wtwo{\iatom}^{(\ilayer)}
      \]

    \column{0.45\textwidth}
      \begin{shibiunote}[Benefit]
        Decouples us from exact knowledge
        of the propagation model through
        the wall -- the network learns it.
      \end{shibiunote}
  \end{columns}
\end{frame}

% ------------------------------------------------------------------
% SLIDE 4 -- LCRPCA: Unrolled Architecture
% ------------------------------------------------------------------
\begin{frame}{LCRPCA: Unrolled Architecture}
  \begin{center}
  \small
  % lcrpca_arch.tex
% LCRPCA architecture: one layer of the unrolled network.
% Two parallel update paths:
%   Path 1 (L path): SVT (singular value thresholding)
%   Path 2 (R path): conv W_{1,k}, soft-threshold S_lambda, conv W_{2,k}
% Shared residual computation.
%
% Colors: theme names
% Bounding box: approx 12cm x 6cm

\begin{tikzpicture}[
    % Operation box
    opbox/.style = {
      draw=shibuiaccent,
      fill=shibuicodebg,
      rounded corners=3pt,
      thick,
      text=shibuitext,
      font=\scriptsize,
      align=center,
      minimum width=2.0cm,
      minimum height=0.9cm,
      inner sep=4pt
    },
    % Data node (small circle)
    datanode/.style = {
      draw=shibuitext,
      fill=shibuicodebg!70!black,
      circle,
      minimum size=0.45cm,
      inner sep=0pt,
      font=\scriptsize,
      text=shibuitext
    },
    % Addition node
    addnode/.style = {
      draw=shibuitext,
      fill=shibuicodebg,
      circle,
      minimum size=0.5cm,
      inner sep=0pt,
      font=\scriptsize\bfseries,
      text=shibuitext
    },
    arr/.style = {->, draw=shibuitext, thick,
                  shorten >=2pt, shorten <=2pt},
    arrthin/.style = {->, draw=shibuitext, thin, dashed,
                      shorten >=2pt, shorten <=2pt},
    lbl/.style = {font=\small, text=shibuitext}
  ]

  % -------- Input Y --------
  \node[datanode] (Y) at (0, 3.0) {$\Mimage$};

  % -------- Residual computation --------
  % Residual R = Y - L - sum_k W_{1,k} * R_k
  \node[opbox, minimum width=3.0cm] (resid) at (3.0, 3.0)
    {Residual\\[2pt]
     $\Mimage - \Mwall - \sum_\iatom \Wone{\iatom} \Fconv \Mtarget_\iatom$};

  \draw[arr] (Y.east) -- (resid.west);

  % -------- Split into two paths --------
  % Branch point
  \coordinate (split) at (5.5, 3.0);
  \draw[arr] (resid.east) -- (split);

  % ==== L PATH (top) ====
  % Step: L += w_0 * residual, then SVT
  \node[opbox] (Lstep) at (5.5, 4.5)
    {$\Mwall + \lparamw{t} \cdot \text{res}$};
  \node[opbox, draw=shibuiplot3, text=shibuiplot3] (Lsvt) at (8.0, 4.5)
    {$\pxSvT{\lparammu{t}}(\cdot)$\\[2pt]
     \textcolor{shibuiplot3}{SVT}};
  \node[datanode] (Lout) at (10.5, 4.5) {$\Mwall^{(t)}$};

  \draw[arr] (split) |- (Lstep.west);
  \draw[arr] (Lstep.east) -- (Lsvt.west);
  \draw[arr] (Lsvt.east) -- (Lout.west);

  % L path label
  \node[lbl, text=shibuiplot3, font=\tiny\bfseries, above=2pt]
    at (8.0, 5.15) {low-rank path};

  % ==== R PATH (bottom) ====
  % Step: R_k += W_{2,k} * residual, then soft-threshold
  \node[opbox] (Rstep) at (5.5, 1.5)
    {$\Mtarget_\iatom + \Wtwo{\iatom} \Fconv \text{res}$};
  \node[opbox, draw=shibuiplot2, text=shibuiplot2] (Rsft) at (8.0, 1.5)
    {$\pxSoftT{\lparamlam{t}}(\cdot)$\\[2pt]
     \textcolor{shibuiplot2}{soft-thresh}};
  \node[datanode] (Rout) at (10.5, 1.5) {$\Mtarget^{(t)}$};

  \draw[arr] (split) |- (Rstep.west);
  \draw[arr] (Rstep.east) -- (Rsft.west);
  \draw[arr] (Rsft.east) -- (Rout.west);

  % R path label
  \node[lbl, text=shibuiplot2, font=\tiny\bfseries, below=2pt]
    at (8.0, 0.85) {sparse path};

  % ==== Feedback to residual (next layer) ====
  \draw[arrthin, rounded corners=4pt]
    (Lout.east) -- (11.5, 4.5) -- (11.5, 3.0) -- (3.0, 3.0 - 0.45)
    node[lbl, font=\tiny, midway, right=2pt, rotate=90,
         anchor=south] {feedback to layer $t+1$};
  \draw[arrthin, rounded corners=4pt]
    (Rout.east) -- (11.5, 1.5);

  % -------- Layer bracket --------
  \draw[decorate, decoration={brace, amplitude=5pt}, shibuitext]
    (-0.3, 5.55) -- (-0.3, 0.5)
    node[lbl, font=\scriptsize, midway, left=8pt, rotate=90,
         anchor=south] {Layer $t$};

  % -------- Learnable parameters annotation --------
  \node[lbl, font=\tiny, align=left, text=shibuiaccent]
    at (6.2, -0.3)
    {Learnable:
     $\lparamw{t}$,
     $\lparammu{t}$,
     $\lparamlam{t}$,
     $\{\Wone{\iatom}^{(t)}, \Wtwo{\iatom}^{(t)}\}_{\iatom=1}^{\nbatoms}$};

\end{tikzpicture}

  \end{center}

  \vspace{0.15em}
  \begin{shibuiframe}{ }
    \small\centering
    One layer: parallel \emphtext{low-rank} ($\Mwall$) and \emphtext{sparse} ($\Mtarget$) paths.
    Learnable: $\lparamlam{\ilayer}, \lparammu{\ilayer}, \lparamw{\ilayer}$, filters.
  \end{shibuiframe}

  \vspace{0.3em}
  \begin{columns}[T]
    \column{0.33\textwidth}
      \small\centering
      $\nlayer = 6$ layers

    \column{0.34\textwidth}
      \small\centering
      $K = 64$ filters

    \column{0.33\textwidth}
      \small\centering
      $7{\times}7 \to 3{\times}3$
  \end{columns}
\end{frame}

% ------------------------------------------------------------------
% SLIDE 5 -- LCRPCA: Architecture Details  [skip?]
% ------------------------------------------------------------------
\begin{frame}{LCRPCA: Architecture Details \skipbadge{}}

  \begin{shibuiframe}{Per-layer updates}
    \small
    $\Mwall^{(\ilayer)} = \pxSvT{\lparammu{\ilayer}}\!\left(\Mwall^{(\ilayer-1)} + \lparamw{\ilayer}(\Mtwri - \text{residual})\right)$

    $\Mtarget^{(\ilayer)} = \pxSoftT{\lparamlam{\ilayer}}\!\left(\Mtarget^{(\ilayer-1)} + \Wtwo{\ilayer}\Fconv(\text{residual})\right)$
  \end{shibuiframe}

  \vspace{0.3em}
  \begin{columns}[T]
    \column{0.50\textwidth}
      \small
      \textbf{Parameters:}

      \vspace{0.2em}
      \renewcommand{\arraystretch}{1.1}
      \begin{tabular}{lc}
        \toprule
        \textbf{Model} & $\nparams$ \\
        \midrule
        LCRPCA & $21{,}656$ \\
        U-Net & $8{,}6\mathrm{M}$ \\
        \midrule
        Ratio & \textbf{400$\times$} \\
        \bottomrule
      \end{tabular}

    \column{0.50\textwidth}
      \small
      \textbf{Training loss:}

      \vspace{0.2em}
      $\ell_2$ (reconstruction) + $\mathcal{H}$ (cross-entropy) + Dice
  \end{columns}
\end{frame}

% ------------------------------------------------------------------
% SLIDE 6 -- LCRPCA: Results with Scarce Data
% ------------------------------------------------------------------
\begin{frame}{LCRPCA: Results with Scarce Data}

  \begin{columns}[T]
    \column{0.50\textwidth}
      \centering
      \figplaceholder[0.92\linewidth]{PR-AUC vs.\ training data
        fraction \\ LCRPCA vs.\ U-Net vs.\ CoroNA}

    \column{0.50\textwidth}
      \centering
      \figplaceholder[0.92\linewidth]{TWRI image reconstruction \\
        LCRPCA (10\% data) vs.\ U-Net (10\% data)}
  \end{columns}

  \vspace{0.5em}
  \begin{shibuiframe}{Key result (10\% training data)}
    \centering
    LCRPCA \emphtext{PR-AUC $= 0.85$}
    \quad vs.\quad
    U-Net $= 0.80$
    \quad vs.\quad
    CoroNA $= 0.07$
  \end{shibuiframe}

  \vspace{0.3em}
  \citebadge{EUSIPCO 2024}
\end{frame}

% ------------------------------------------------------------------
% SLIDE 7 -- Why Fewer Parameters Help
% ------------------------------------------------------------------
\begin{frame}{Why Fewer Parameters Help}

  \begin{columns}[T]
    % ------- LEFT: explanation -------
    \column{0.55\textwidth}
      \vspace{0.3em}
      \begin{itemize}
        \item \emphtext{Standard nets}: learn everything, no prior
        \pause
        \item \emphtext{LCRPCA}: algorithm structure provides bias
        \pause
        \item Structure: low-rank + sparsity proxes + conv layers
        \pause
        \item Fewer parameters $\Rightarrow$ less data needed
        \pause
        \item Dictionary learned from data
      \end{itemize}

    % ------- RIGHT: key inequality -------
    \column{0.45\textwidth}
      \vspace{0.3em}
      \begin{shibuiframe}{Parameter count}
        \[
          \nparams(\mathrm{LCRPCA})
            = 21{,}656
            \ll
          \nparams(\mathrm{U\text{-}Net})
            = 8{,}650{,}474
        \]
      \end{shibuiframe}

      \vspace{0.6em}
      \begin{shibuitip}
        LCRPCA generalizes better in the
        \emphtext{scarce-data regime} precisely
        because its inductive bias reduces
        the effective hypothesis space.
      \end{shibuitip}
  \end{columns}
\end{frame}

% ------------------------------------------------------------------
% SLIDE 8 -- Comparison: Classical vs.\ Learned
% ------------------------------------------------------------------
\begin{frame}{Comparison: Classical vs.\ Learned}

  \vspace{0.1em}
  \small
  \renewcommand{\arraystretch}{1.25}
  \begin{center}
    \begin{tabular}{lcccc}
      \toprule
      \textbf{Method} & \textbf{Params} & \textbf{10\%} & \textbf{100\%} & \textbf{Interp.} \\
      \midrule
      PGD (KRPCA) & 0 & 0.62 & 0.70 & Physics \\
      U-Net & 8.6M & 0.80 & 0.87 & Black box \\
      CoroNA & -- & 0.07 & -- & Partial \\
      \textbf{LCRPCA} & \textbf{21k} & \textbf{0.85} & \textbf{0.89} & Unrolled \\
      \bottomrule
    \end{tabular}
  \end{center}

  \vspace{0.4em}
  \begin{shibuitip}[Take-away]
    Best of both: \emphtext{interpretability} (algorithm) + \emphtext{flexibility} (learning)
  \end{shibuitip}
\end{frame}
